% ---
% id: statistics/binomial-box/rem:selector
% kind: rem
% title: "The role of the constrained outcome set"
% summary: >-
%   Constraining fewer outcomes narrows each interval but leaves more
%   coordinates bounded only by the simplex, so no choice of the constrained
%   outcomes dominates. Completeness holds for every choice, and the
%   unconstrained outcomes keep the bounds 0 and 1, so the box always has two
%   rows per outcome.
% prominence: standard
% status: accepted
% uses: [def:construction]
% mentions: [thm:completeness-bound, req:request-defaults]
% assumes: []
% refs: []
% provenance: newly-derived
% lean:
% ---

Which outcomes are constrained is part of the input; the default, all of
them, is Requirement~\ref{req:request-defaults}.

\begin{remark}\label{rem:selector}
Constraining fewer outcomes concentrates the budget on those that are
constrained, which narrows each interval; it also weakens the set, because the
unconstrained coordinates are then bounded only by the simplex. Both effects
are real and pull in opposite directions, so no choice of $\Chatcon$
dominates. Completeness holds for whatever $\Chatcon$ is supplied
(Theorem~\ref{thm:completeness-bound}), and the unconstrained coordinates receive the vacuous
bounds $[0,1]$ of \eqref{eq:boxbounds} rather than being omitted from $B$ ---
so the row count $m_{\mathrm{row}}=2|\Chatset|$ does not depend on $\Chatcon$
and a consumer never has to infer a missing row.
\end{remark}
